% SOURCE_AUTHORITY: sga5_fr_workpass.tex, lines 11627--11667 (LF-normalized unit).
% SOURCE_SHA256: 791F4EFFC5E02832D5D77ED03518C8156D6F07E4C8238B03545DB93D883FBB28
Para demonstrar $\delta)$ são utilizadas as sequências exatas:
\[
0\longrightarrow Z^i\longrightarrow E^i\longrightarrow B^{i+1}\longrightarrow 0,
\]
\[
0\longrightarrow B^i\longrightarrow Z^i\longrightarrow H^i(E)\longrightarrow 0.
\]
A demonstração de $\varepsilon)$ é feita por indução sobre $n$, utilizando a sequência exata
\[
0\longrightarrow F_{n-1}\longrightarrow F\longrightarrow F/F_{n-1}\longrightarrow 0.
\]
Para demonstrar $\varphi)$, sejam $E\in\Ob\cD$ e
\[
F=\bigoplus_{n\in\mathbb Z^+}E_n,\qquad E_n=E,\quad n\in\mathbb Z^+.
\]
Então $E\oplus F\simeq F$ e, aplicando $\gamma)$, temos
\[
f(F)=f(F)+f(E),
\]
Por isso, $f(E)=0$.

Se considerarmos, para todo grupo abeliano $G$, o grupo abeliano $A_{\cD}(G)$ das funções aditivas de $\cD$ em $G$, obteremos evidentemente um funtor covariante
\[
A_{\cD}:\Ab\longrightarrow\Ab.
\]
O funtor $A_{\cD}$ é representável. Um par de representação $(K_{\cD}(\cC),\varphi)$ pode ser obtido da seguinte forma. Denotando por $L$ o grupo abeliano livre gerado pelos objetos de $\cD$, seja $R$ o subgrupo de $L$ gerado pelos elementos da forma $E-E'-E''$ para cada sequência exata
\[
0\longrightarrow E'\longrightarrow E\longrightarrow E''\longrightarrow 0
\]
de $\cC$. Põe-se
\[
K_{\cD}(\cC)=L/R.
\]
A aplicação $\varphi$ é obtida compondo as duas aplicações canônicas
\[
\Ob\cD\longrightarrow L\longrightarrow L/R.
\]
Quando não há perigo de confusão, a notação pode ser abreviada:
\[
K_{\cD}(\cC)=K(\cD).
\]
